\section{Methods}

\subsection{Data and splits}

We use HAM10000, which contains 10,015 dermoscopic images from 7,470 lesions across seven diagnostic classes \citep{tschandl2018ham10000}. The task is binary: melanoma versus all other diagnoses. The dataset contains 1,113 melanoma images, 11.1\% of the total.

We use a lesion-aware split so that images of the same lesion never appear in more than one split. The split is 70\% training, 15\% validation, and 15\% test. The test set contains 1,527 images and 186 melanomas. The split is generated with seed 42, released as \texttt{results/split\_ham10000\_lesion\_seed42.csv}, and reused by every run.

\subsection{Preprocessing and augmentation}

Images are resized to 160 pixels for the Stage 1 screen and 224 pixels for Stages 2 and 3, and normalized with ImageNet statistics. The standard augmentation pipeline applies horizontal and vertical flips, random 90-degree rotations, and mild color jitter. The strong pipeline increases the color jitter and adds affine transformations with scale, translation, and rotation, as a simulation of acquisition shift.

\subsection{Architectures}

We compare ResNet-50, EfficientNet-B0, and ViT-S/16. All models use ImageNet pretrained weights and a two-class output. The ViT uses dynamic image sizing so that it accepts the 224-pixel input.

\subsection{Training}

All models use AdamW with weight decay $1\times10^{-4}$, a cosine schedule, a batch size of 64, at most 15 epochs, gradient clipping at a norm of 1.0, and mixed precision on GPU. Early stopping uses a patience of 5 epochs on the checkpoint selection metric, and the checkpoint with the best value of that metric is kept. Focal loss uses $\gamma=2$ and $\alpha=0.75$ on the melanoma class in every focal configuration. The learning rate and the selection metric depend on the protocol (Section~\ref{sec:protocols}).

\subsection{Protocols and conditions}
\label{sec:protocols}

\paragraph{Stage 1.} A screen on EfficientNet-B0 at 160 pixels with eight configurations: CE with a random sampler, CE with class-balanced sampling, weighted CE, focal loss with a random sampler, focal loss with class-balanced sampling, focal loss with $\gamma=3$ and class-balanced sampling, and two strong-augmentation variants. Each configuration uses three seeds, for 24 runs. Stage 1 uses the original protocol below.

\paragraph{Stage 2, original protocol.} Three pre-declared conditions across the three architectures: the CE baseline (cross-entropy, random sampler, standard augmentation), focal loss with class-balanced sampling (Focal+CB), and the Stage 1 winner, CE with strong augmentation (CE+strong). All architectures share a learning rate of $3\times10^{-4}$. The checkpoint is selected, and early stopping is applied, on validation balanced accuracy at the 0.5 probability threshold. Each condition uses three seeds, for 27 runs.

\paragraph{Audit of the original protocol.} Two features of the original protocol could favor one architecture-condition combination over another. First, balanced accuracy at 0.5 depends on where the scores fall relative to 0.5, and class-balanced sampling and the focal $\alpha$ move the scores toward melanoma. The rule therefore selects checkpoints differently by condition. In Stage 2 the focal CNN checkpoints came earlier (median epoch 4) than the CE CNN checkpoints (median 8), while the CE ViT checkpoints came very early (epochs 3, 1, and 5) and the focal ViT checkpoints later (median 8). Second, a learning rate shared across architectures may suit some architectures better than others. We designed Stage 3 to remove both features. Stage 3 was designed after the Stage 2 test results were known, so it is a targeted re-test of a specific result rather than a blind replication. All Stage 3 settings below were fixed before any Stage 3 run was trained.

\paragraph{Learning rate selection.} For each architecture we trained the CE baseline with seed 1 at learning rates $3\times10^{-5}$, $1\times10^{-4}$, and $3\times10^{-4}$, with checkpoint selection by validation AUC, and kept the learning rate with the highest validation AUC. The selected values are $3\times10^{-4}$ for ResNet-50 (validation AUC 0.889, against 0.885 and 0.853), $3\times10^{-4}$ for EfficientNet-B0 (0.907, against 0.886 and 0.863), and $3\times10^{-5}$ for ViT-S/16 (0.906, against 0.894 and 0.893). Supplementary Table~S1 lists all nine runs. The test split was not used for this choice.

\paragraph{Stage 3, revised protocol.} Each architecture uses its selected learning rate. The checkpoint is selected, and early stopping is applied, on validation AUC, which does not depend on a threshold. Stage 3 keeps the three Stage 2 conditions and adds two ablations that separate the components of Focal+CB: focal loss with a random sampler (Focal) and cross-entropy with class-balanced sampling (CE+CB). Each of the five conditions uses five seeds, for 75 runs. All other settings match Stage 2.

Stage 3 is the primary benchmark. The fairness, calibration, explanation, and external analyses use the Stage 3 models. Stage 2 is reported to show how the protocol changes the conclusion.

\subsection{Metrics and threshold}

We report sensitivity, specificity, balanced accuracy, AUC, and sensitivity at 95\% specificity. The operating threshold is selected on the validation split at 95\% specificity, frozen, and applied once to the test split. Each run writes per-sample test predictions so that the same scores are reused for the fairness and XAI analyses.

\subsection{Statistical analysis}

The seed is the experimental unit for a training intervention. We report the mean across seeds with a 95\% t interval, and release the per-seed values. Image-level bootstrap intervals are not used as confidence intervals for a seed-level mean. Paired changes compare a condition with the CE baseline of the same architecture and seed on the same test images. The architecture interaction is a difference in differences: the condition-minus-CE change of one architecture minus that of another, computed per seed. AUC differences between architectures are paired per seed in the same way. We also compute sensitivity at a common 95\% test specificity for every run. The direction count reports how many seeds share the sign of the mean. The intervals are not corrected for the number of comparisons, so an isolated interval that excludes zero is weak evidence.

\subsection{Fairness analysis}

HAM10000 has no Fitzpatrick labels. We therefore compute an image-based lightness proxy from non-lesional skin. Lesion masks are the HAM10000 segmentation masks released by Tschandl et al.\ \citep{tschandl2020human}. The image is converted to CIELAB, the lesion mask is dilated by 11 pixels and excluded, and very dark pixels ($L^*<10$) and very bright pixels ($L^*>95$) are removed. The median $L^*$ of the remaining skin pixels is the proxy.

We group images into light ($L^*>70$), medium ($55<L^*\leq70$), and dark ($L^*\leq55$). These thresholds are specific to the HAM10000 lightness distribution. They are image-based strata, not clinical Fitzpatrick labels.

We also computed the Individual Typology Angle (ITA), $\mathrm{ITA}=\arctan((L^*-50)/b^*)\times180/\pi$. The ITA was unstable on 37\% of images because the $b^*$ channel was near zero or negative after dermatoscope color casts. The $L^*$ proxy is therefore the primary analysis, and the ITA is not used for group comparisons.

For each run, we compute per-group sensitivity, specificity, false positive rate, false negative rate, and AUC, and report the mean and standard deviation across seeds.

Lesion contrast is computed from the same masks. The lesion color is the median CIELAB value inside the mask, and the surrounding skin is the median inside an 11-pixel ring after excluding $L^*<10$ and $L^*>95$. We define $\Delta L^*$ as skin minus lesion, so a positive value means the lesion is darker than the surrounding skin.

The main analysis uses the pre-declared cutoff $L^*\leq55$. We repeat the group analysis at $L^*\leq58$ and $L^*\leq60$ to test whether the results depend on the cutoff, and compare dark-stratum AUC with the light and medium strata by bootstrap for every CE-baseline run. We also fit adjusted logistic models of false positive and false negative on perilesional $L^*$, $\Delta L^*$, log lesion area, age, sex, anatomical site, and acquisition source, using the CE-baseline predictions of each architecture. HAM10000 combines four source collections that differ in clinic, dermatoscope, and melanoma prevalence (0.6\% to 20\%), and image lightness depends on the device, so source is a potential confounder of the lightness effect. Diagnosis confirmation type is not included because it depends on the label. Continuous predictors are standardized; one standard deviation is 7.2 $L^*$ units and 10.3 contrast units in the test set. Standard errors come from a cluster bootstrap over lesions with 300 replicates, so repeated images and seeds of the same lesion are not treated as independent. The $p$-values are descriptive and are not corrected for multiplicity.

\subsection{Calibration}

We compute expected calibration error (ECE) with 10 equal-width bins and the Brier score on the test split. Class-balanced sampling and the focal $\alpha$ train the model at an effective melanoma prior above the 11\% prevalence. We therefore also report ECE after two monotone corrections: a prior-shift correction that adds the log ratio of the prevalence odds to the effective training odds to each logit, followed by a temperature fitted on the validation split by minimizing the negative log-likelihood. Neither correction changes the AUC or the validation-selected operating point.

\subsection{External validation}

External validation uses HIBA \citep{hiba2023}, a dermoscopic dataset from Argentina with Fitzpatrick labels. We evaluate the CE, Focal+CB, and CE+strong models of every architecture and seed. We apply the frozen HAM10000 validation threshold and do not refit on HIBA. As a secondary comparison, we also report sensitivity at 95\% specificity with a threshold fitted on HIBA. This oracle threshold is an upper bound, not a deployable result. PROVe-AI \citep{proveai2023} is a possible extension but was not evaluated here. A clinical dark-skin dataset, such as DDI \citep{daneshjou2022ddi} or PASSION \citep{passion2024}, is left for future work because of the imaging modality difference.

\subsection{XAI analysis}

For the CE-baseline run of each architecture and seed, we compute Grad-CAM heatmaps on every test image: 644 light, 840 medium, and 43 dark images per seed. Grad-CAM targets the true class for every image, including false negatives, so the heatmap describes evidence for the true class rather than the predicted class. We measure:

\begin{itemize}
  \item Attribution mass inside the lesion: $M=\sum_{i\in\mathrm{lesion}}h_i/\sum_i h_i$, where $h_i$ is the normalized heatmap value at pixel $i$.
  \item Pointing game accuracy: the fraction of images where the maximum heatmap value falls inside the lesion mask.
  \item Attribution IoU: the intersection over union between the thresholded heatmap and the lesion mask.
\end{itemize}

We then compute the Spearman correlation between attribution mass and error status, and between attribution mass and lightness group. The correlations pool the five seeds, so the $p$-values treat repeated images as independent and are descriptive.

We also implemented deletion and insertion curves, but in the original models the target-class probability stayed near $1.0$ for both the original and the perturbed images, so the curves did not separate heatmaps. We therefore report localization metrics only.

\subsection{Reproducibility}

Every run writes its configuration, git commit, seed, device, package versions, metrics, and per-sample predictions. The split is cached and reused. Training used Kaggle GPUs. The calibration, XAI, and external evaluations used a local Apple Silicon GPU. The code, configurations, and table generation scripts are available in the project repository.
